%% LyX 2.5.1 created this file.  For more info, see https://www.lyx.org/.
%% Do not edit unless you really know what you are doing.
\documentclass[10pt,english,10pt]{beamer}
\usepackage{lmodern}
\usepackage[T1]{fontenc}
\usepackage[utf8]{inputenc}
\setcounter{tocdepth}{1}
\usepackage{babel}
\usepackage{amsbsy}
\usepackage{amstext}
\usepackage{amssymb}
\usepackage{graphicx}
\usepackage[authoryear]{natbib}
\ifx\hypersetup\undefined
  \AtBeginDocument{%
    \hypersetup{}
  }
\else
  \hypersetup{}
\fi

\makeatletter
%%%%%%%%%%%%%%%%%%%%%%%%%%%%%% Textclass specific LaTeX commands.
% this default might be overridden by plain title style
\newcommand\makebeamertitle{\frame{\maketitle}}%
% (ERT) argument for the TOC
\AtBeginDocument{%
  \let\origtableofcontents=\tableofcontents
  \def\tableofcontents{\@ifnextchar[{\origtableofcontents}{\gobbletableofcontents}}
  \def\gobbletableofcontents#1{\origtableofcontents}
}

%%%%%%%%%%%%%%%%%%%%%%%%%%%%%% User specified LaTeX commands.



\usepackage{tikz}
\usetikzlibrary{positioning}
\usetikzlibrary{overlay-beamer-styles}
\usepackage{appendixnumberbeamer}

\usepackage{graphicx}
\usepackage{subfig}

\usetheme[progressbar=frametitle,block=fill,subsectionpage=progressbar]{metropolis}

% margin
\setbeamersize{text margin right=1.5cm}

% colors
\definecolor{DarkRed}{rgb}{0.7,0,0}
%\colorlet{DarkRed}{red!70!black}
\setbeamercolor{normal text}{fg=black}
\setbeamercolor{alerted text}{fg=DarkRed}
\setbeamercolor{progress bar}{fg=DarkRed}
\setbeamercolor{button}{bg=DarkRed}

% width of seperators
\makeatletter
\setlength{\metropolis@titleseparator@linewidth}{1pt}
\setlength{\metropolis@progressonsectionpage@linewidth}{1pt}
\setlength{\metropolis@progressinheadfoot@linewidth}{1pt}
\makeatother

% new alert block
\newlength\origleftmargini
\setlength\origleftmargini\leftmargini
\setbeamertemplate{itemize/enumerate body begin}{\setlength{\leftmargini}{4mm}}
\let\oldalertblock\alertblock
\let\oldendalertblock\endalertblock
\def\alertblock{\begingroup \setbeamertemplate{itemize/enumerate body begin}{\setlength{\leftmargini}{\origleftmargini}} \oldalertblock}
\def\endalertblock{\oldendalertblock \endgroup}
\setbeamertemplate{mini frame}{}
\setbeamertemplate{mini frame in current section}{}
\setbeamertemplate{mini frame in current subsection}{}
\setbeamercolor{section in head/foot}{fg=normal text.bg, bg=structure.fg}
\setbeamercolor{subsection in head/foot}{fg=normal text.bg, bg=structure.fg}

% footer
\makeatletter
\setbeamertemplate{footline}{%
    \begin{beamercolorbox}[colsep=1.5pt]{upper separation line head}
    \end{beamercolorbox}
    \begin{beamercolorbox}{section in head/foot}
      \vskip1pt\insertsectionnavigationhorizontal{\paperwidth}{}{\hskip0pt plus1filll \insertframenumber{} / \inserttotalframenumber \hskip2pt}\vskip3pt%
    \end{beamercolorbox}%
    \begin{beamercolorbox}[colsep=1.5pt]{lower separation line head}
    \end{beamercolorbox}
}
\makeatother

% toc
\setbeamertemplate{section in toc}{\hspace*{1em}\inserttocsectionnumber.~\inserttocsection\par}
\setbeamertemplate{subsection in toc}{\hspace*{2em}\inserttocsectionnumber.\inserttocsubsectionnumber.~\inserttocsubsection\par}
% Added by lyx2lyx
\setlength{\parskip}{\smallskipamount}
\setlength{\parindent}{0pt}

\makeatother

\begin{document}
\title{4. Stationary Equilibrium\vspace{-2mm}}
\subtitle{Adv. Macro: Heterogeneous Agent Models}
\author{Raphaël Huleux}
\date{2026}

{
\setbeamertemplate{footline}{}
\begin{frame}

\maketitle

\end{frame}
}

\addtocounter{framenumber}{-1}
\begin{frame}{\only<1|handout:0>{So far: the household block}\only<2|handout:0>{Today: prices come from firms}\only<3>{Today: markets close the model}}
\vspace{2mm}

\begin{center}
\begin{tikzpicture}[
block/.style={draw,thick,rounded corners,minimum width=3.1cm,minimum height=1cm,align=center},
lbl/.style={font=\scriptsize},
new/.style={draw=DarkRed,text=DarkRed}]

\node[block] (hh) at (0,0) {Households};
\node[block,visible on=<2->,alt=<2>{new,fill=DarkRed!10}{}] (firm) at (6.8,0) {Firms};
\node[block,visible on=<3->,new,fill=DarkRed!10] (asset) at (3.4,1.85) {Asset market\\$A_{t}=K_{t}$};
\node[block,visible on=<3->,new,fill=DarkRed!10] (goods) at (3.4,-1.85) {Goods market\\$Y_{t}=C_{t}+I_{t}$};

\draw[->,thick,alt=<2>{new}{}] ([yshift=2.5mm]firm.west) -- node[lbl,above]{\alt<1>{$r_{t}$ given}{$r_{t}=F_{K}(K_{t-1},L_{t})-\delta$}} ([yshift=2.5mm]hh.east);
\draw[->,thick,alt=<2>{new}{}] ([yshift=-2.5mm]firm.west) -- node[lbl,below]{\alt<1>{$w_{t}$ given}{$w_{t}=F_{L}(K_{t-1},L_{t})$}} ([yshift=-2.5mm]hh.east);

\draw[->,thick] (hh.north) |- node[lbl,pos=0.25,left]{$A_{t}$} (asset.west);
\draw[->,thick] (hh.south) |- node[lbl,pos=0.25,left]{$C_{t}$} (goods.west);
\draw[->,thick,visible on=<2->,alt=<2>{new}{}] (asset.east) -| node[lbl,pos=0.75,right]{$K_{t-1}$} (firm.north);
\draw[->,thick,visible on=<2->,alt=<2>{new}{}] (firm.south) |- node[lbl,pos=0.25,right]{$Y_{t}$, $I_{t}$} (goods.east);

\end{tikzpicture}
\end{center}

\vspace{1mm}

\begin{overlayarea}{\textwidth}{1.85cm}
\only<1|handout:0>{\textbf{Partial equilibrium: }households take $r_{t}$ and $w_{t}$ as given and choose $c_{t}$ and $a_{t}$.}
\only<2|handout:0>{\textbf{Firms: }$r_{t}$ and $w_{t}$ are now determined by the capital stock.

\vspace{2mm}
But who decides $K_{t}$?}
\only<3>{\textbf{General equilibrium: }capital is what households save, $K_{t}=A_{t}$.

\vspace{2mm}
Prices must be such that both markets clear: a fixed point.}
\end{overlayarea}
\end{frame}
%
\begin{frame}{Recap of the two last classes}

What did we learn so far?

\onslide<+->
\begin{itemize}
\item <+->The buffer-stock model captures some facts about consumption
and MPCs
\item <+->We can solve it with dynamic programming:
\begin{itemize}
\item <+->The Value Function Iteration algorithm (slow)
\item <+->The Endogenous Grid Method (fast)
\item <+->The goal is to obtain the policy functions $c(a,z)$, $a'(a,z)$
\end{itemize}
\item <+->How to simulate a distribution of households using
\begin{itemize}
\item <+->The Monte-Carlo method (slow and imprecise)
\item <+->The histogram method (fast and precise)
\end{itemize}
\end{itemize}
\onslide<+->Today and next week: use those methods to solve the steady-state
of a simple heterogeneous-agent model.
\end{frame}
%

\section{Introduction}
\begin{frame}{Introduction}
\vspace{-2mm}
\begin{itemize}
\item <+->\textbf{Last time: }
\begin{enumerate}
\item Partial equilibrium
\item No interactions
\end{enumerate}
\item <+->\textbf{Today:} Interaction through markets
\item <+->\textbf{Model:} Heterogeneous Agent Neo-Classical (HANC) model
\item <+->\textbf{Equilibrium-concept:} \emph{Stationary equilibrium}
\begin{enumerate}
\item What determines income and wealth inequality\emph{ in the long run}?
\item What determines the real interest rate \emph{in the long run}?
\end{enumerate}
\item <+->\textbf{Code:} Based on the \textcolor{DarkRed}{\href{https://github.com/NumEconCopenhagen/GEModelTools}{GEModelTools}}
package
\begin{enumerate}
\item Is in active development
\item You can help to improve interface, find bugs and features
\end{enumerate}
\textbf{Documentation: }See \textcolor{DarkRed}{\href{https://github.com/NumEconCopenhagen/GEModelToolsNotebooks}{GEModelToolsNotebooks}}
\begin{itemize}
\item Many examples in repo, so look if you have issues
\end{itemize}
\item <+->\textbf{Literature:} Aiyagari (1994)
\end{itemize}
\end{frame}
%
\begin{frame}{Outline of this lecture}
\begin{enumerate}
\item Recap of the Ramsey (Neo-Classical) model
\item Overview of the Heterogeneous-Agent Neo-Classical model (HANC)
\item How to compute the stationary equilibrium
\item Some economic properties of the HANC stationary equilibrium
\end{enumerate}
\end{frame}
%

\section{Ramsey-recap}
\begin{frame}{The Ramsey model}
\begin{itemize}
\item <+->We will study the \emph{stationary equilibrium }in\emph{ }the\emph{
}Heterogeneous Agent Neo-Classical (HANC) model
\item <+->Merges two well known models in the literature:
\begin{itemize}
\item <+->Standard Ramsey–Cass–Koopman model (\emph{NC})
\begin{itemize}
\item What do we mean by Neo-classical?
\end{itemize}
\item <+->One-asset Buffer-stock model (\emph{HA})
\end{itemize}
\item <+->Went through the Buffer-stock model over the last two lectures
\item <+->\textbf{Now}: Recap of the Ramsey model
\end{itemize}
\end{frame}
%
\begin{frame}{Ramsey: Firms}
\begin{itemize}
\item <+->\textbf{Production function: $Y_{t}=F(\Gamma_{t},K_{t-1},L_{t})$ }{\footnotesize{[}capital
chosen in $t-1$ is used for production at $t${]}}{\footnotesize\par}

where $\Gamma_{t}$ is technology
\item <+->\textbf{Profits:} $\Pi_{t}=Y_{t}-w_{t}L_{t}-r^{K}_{t}K_{t-1}$
\item <+->\textbf{Profit maximization: $\max_{K_{t-1},L_{t}}\Pi_{t}$}\vspace{1mm}
\begin{enumerate}
\item Rental rate: $\frac{\partial\Pi_{t}}{\partial K_{t-1}}=0\Leftrightarrow$$r^{K}_{t}=F_{K}(\Gamma_{t},K_{t-1},L_{t})$\vspace{1mm}
\item Real wage: $\frac{\partial\Pi_{t}}{\partial L_{t}}=0\Leftrightarrow w_{t}=F_{L}(\Gamma_{t},K_{t-1},L_{t})$\vspace{1mm}
\end{enumerate}
\vspace{1mm}With CRS we get zero profits: $\Pi_{t}=0\Rightarrow$\vspace{2mm}

$Y_{t}=w_{t}L_{t}+r^{K}_{t}K_{t-1}$ {\footnotesize{[}functional income
distribution{]}}{\footnotesize\par}
\end{itemize}
\end{frame}
%
\begin{frame}{Ramsey: Zero-profit mutual fund}
\begin{itemize}
\item <+->Introduce \textbf{mutual fund}
\begin{itemize}
\item Takes savings $A_{t-1}$ from households and invest them in available
assets
\item In the Ramsey model: Only capital $K_{t-1}$ but could also include
gov. bonds, firm equity etc.
\item Receive income from firms and redistribute it to households
\end{itemize}
\item <+->\textbf{Capital depreciate }with rate $\delta\in(0,1)$,\textbf{
\[
K_{t}=(1-\delta)K_{t-1}+I_{t}
\]
}
\item <+->\textbf{Deposits }(from households), $A_{t-1}$:\textbf{ }The
rate of return is
\[
r_{t}=r^{K}_{t}-\delta
\]
\item <+->\textbf{Balance sheet:
\[
A_{t-1}=K_{t-1}
\]
}
\end{itemize}
\end{frame}
%
\begin{frame}{Ramsey: Households}
\begin{itemize}
\item \textbf{Utility maximization:}\vspace{-2mm}
\begin{align*}
v_{0}(A^{hh}_{-1}) & =\max_{\{C^{hh}_{t}\}^{\infty}_{t=0}}\sum^{\infty}_{t=0}\beta^{t}u(C^{hh}_{t})\\
 & \text{s.t.}\\
C^{hh}_{t}+A^{hh}_{t} & =(1+r_{t})A^{hh}_{t-1}+w_{t}L^{hh}_{t}
\end{align*}

Exogenous labor supply: $L^{hh}_{t}=1$
\item \textbf{Euler-equation }(implied by Lagrangian):\textbf{
\[
u^{\prime}(C^{hh}_{t})=\beta(1+r_{t+1})u^{\prime}(C^{hh}_{t+1})
\]
}
\end{itemize}
\end{frame}
%
\begin{frame}{Ramsey: Market Clearing}
\begin{itemize}
\item <+->\textbf{Capital market: $K_{t}=A_{t}=A^{hh}_{t}$}
\item <+->\textbf{Labor market: $L_{t}=L^{hh}_{t}=1$}
\item <+->\textbf{Goods market: $Y_{t}=C^{hh}_{t}+I_{t}$}\vspace{2mm}
\item <+->\textbf{Walras:} Capital and labor market clears $\Rightarrow$
goods market clears. Start from
\begin{align*}
C^{hh}_{t}+A^{hh}_{t} & =(1+r_{t})A^{hh}_{t-1}+w_{t}L^{hh}_{t}\\
\Leftrightarrow C^{hh}_{t}+I_{t} & =\left[(1+r_{t})A^{hh}_{t-1}+w_{t}L^{hh}_{t}-A^{hh}_{t}\right]+(K_{t}-(1-\delta)K_{t-1})\\
 & =\left[(1+r_{t})K_{t-1}+w_{t}L_{t}-K_{t}\right]+(K_{t}-(1-\delta)K_{t-1})\\
 & =r^{K}_{t}K_{t-1}+w_{t}L_{t}\\
 & =Y_{t}
\end{align*}
\item <+->Note: Means that we can check if we have solved the numerical
model correctly by:
\begin{itemize}
\item Impose two of the market clearing conditions
\item Then check the third market clearing condition (should be zero)
\end{itemize}
\end{itemize}
\end{frame}
%
\begin{frame}{Ramsey: Summary}
\begin{itemize}
\item <+->\textbf{Extended form:}\vspace{-1mm}\textbf{
\begin{align*}
r^{K}_{t} & =F_{K}(\Gamma_{t},K_{t-1},L_{t})\\
w_{t} & =F_{L}(\Gamma_{t},K_{t-1},L_{t})\\
r_{t} & =r^{K}_{t}-\delta\\
A_{t} & =K_{t}\\
A^{hh}_{t} & =(1+r_{t})A^{hh}_{t-1}+w_{t}L^{hh}_{t}-C^{hh}_{t}\\
u^{\prime}(C^{hh}_{t}) & =\beta(1+r_{t+1})u^{\prime}(C^{hh}_{t+1})\\
A_{t} & =A^{hh}_{t}\\
L_{t} & =L^{hh}_{t}
\end{align*}
}

\item <+->\textbf{Simplified form:}
\begin{align*}
u^{\prime}(C^{hh}_{t}) & =\beta(1+F_{K}(\Gamma_{t+1},K_{t},1)-\delta)u^{\prime}(C^{hh}_{t+1})\\
K_{t} & =(1-\delta)K_{t-1}+F(\Gamma_{t},K_{t-1},1)-C^{hh}_{t}
\end{align*}

\end{itemize}
\end{frame}
%
\begin{frame}{Ramsey: As an equation system}

Eqs. system with unknowns $\left\{ K_{t},L_{t},r^{K}_{t},w_{t},r_{t},A_{t},A^{hh}_{t},C^{hh}_{t}\right\} ^{\infty}_{t=0}$
and eqs:

\begin{align*}
\left[\begin{array}{c}
r^{K}_{t}-F_{K}(\Gamma_{t},K_{t-1},L_{t})\\
w_{t}-F_{L}(\Gamma_{t},K_{t-1},L_{t})\\
r_{t}-(r^{K}_{t}-\delta)\\
A_{t}-K_{t}\\
A^{hh}_{t}-\left((1+r_{t})A^{hh}_{t-1}+w_{t}L^{hh}_{t}-C^{hh}_{t}\right)\\
u^{\prime}(C^{hh}_{t})-\beta(1+r_{t+1})u^{\prime}(C^{hh}_{t+1})\\
A_{t}-A^{hh}_{t}\\
L_{t}-L^{hh}_{t}\\
\forall t\in\{0,1,\dots\},\text{given }K_{-1}
\end{array}\right] & \boldsymbol{=}\boldsymbol{0}
\end{align*}

\end{frame}
%
\begin{frame}{Ramsey: Steady state}
\begin{itemize}
\item \textbf{Euler-equation} can be solved for $r_{ss}$ and hence $K_{ss}$:
\begin{align*}
u^{\prime}(C_{ss}) & =\beta(1+F_{K}(\Gamma_{ss},K_{ss},1)-\delta)u^{\prime}(C_{ss})\Leftrightarrow\\
F_{K}(\Gamma_{ss},K_{ss},1) & =\frac{1}{\beta}-1+\delta
\end{align*}
\item \textbf{Accumulation equation + goods mkt. clearing} then implies
$C_{ss}$:
\begin{align*}
K_{ss} & =(1-\delta)K_{ss}+F(\Gamma_{ss},K_{ss},1)-C_{ss}\Leftrightarrow\\
C_{ss} & =(1-\delta)K_{ss}+F(\Gamma_{ss},K_{ss},1)-K_{ss}
\end{align*}
\item Important thing to note: $r_{ss}=1/\beta-1$ whatever the level of
$K_{ss}$. The steady-state asset supply of households is \emph{perfectly
elastic}, a horizontal line at $1/\beta-1$!
\end{itemize}
\end{frame}
%

\section{HANC}
\begin{frame}{HANC model overview}
\begin{itemize}
\item <+->\textbf{Model blocks:}
\begin{enumerate}
\item \textbf{Firms:} Rent capital from mutual fund and hire labor from
the households, produce with given technology, and sell output goods
\item \textbf{Zero-profit mutual funds:} Own capital and rent it to firms,
take deposits and pay return to household
\item \textbf{Households:} Face idiosyncratic productivity shocks, supplies
labor exogenously and makes consumption-saving decisions
\item \textbf{Markets:} Perfect competition in labor, goods and capital
markets
\end{enumerate}
\item <+->\textbf{Add-on to Ramsey-Cass-Koopman:} \emph{Heterogeneous households
subject to idiosyncratic shocks $\to$ generate precautionary savings! }
\item <+->\textbf{Other names:}
\begin{enumerate}
\item <+->The Aiyagari-model
\item <+->The Aiyagari-Bewley-Hugget-Imrohoroglu-model
\item <+->The Standard Incomplete Market (SIM) model
\end{enumerate}
\end{itemize}
\end{frame}
%
\begin{frame}{Heterogeneous households}
\begin{itemize}
\item <+->\textbf{Utility maximization} for household $i$:\vspace{-2mm}
\begin{align*}
v_{0}(z_{i0},a_{i,-1}) & =\max_{\{c_{it}\}^{\infty}_{t=0}}\mathbb{E}_{0}\sum^{\infty}_{t=0}\beta^{t}u(c_{it})\\
 & \text{s.t.}\\
\ell_{it} & =z_{it}\\
a_{it} & =(1+r_{t})a_{it-1}+w_{t}\ell_{it}-c_{it}\\
\log z_{it+1} & =\rho_{z}\log z_{it}+\psi_{it+1},\,\,\,\psi_{it}\sim\mathcal{N}(\mu_{\psi},\sigma_{\psi}),\,\,\,\mathbb{E}[z_{it}]=1\\
a_{it} & \geq0
\end{align*}
\item <+->\vspace{-1mm}\textbf{Where does heterogeneity enter?}
\begin{enumerate}
\item <+-|handout:0>\emph{Ex post} due to stochastic productivity, $z_{it}$
\end{enumerate}
\item <+->\textbf{Incomplete markets due to borrowing constraint }\\
{\small (fancy words: partial self-insurance, lack of Arrow-Debreu
securities)}{\small\par}
\end{itemize}
\end{frame}
%
\begin{frame}{Recursive formulation}
\begin{itemize}
\item <1->\textbf{Value function }(at decision)
\begin{align*}
v(z_{it},a_{it-1}) & =\max_{c_{it}}u(c_{it})+\beta\mathbb{E}\left[v(z_{it+1},a_{it})\right]\\
 & \text{s.t.}\\
\ell_{it} & =z_{it}\\
a_{it} & =(1+r_{t})a_{it-1}+w_{t}\ell_{it}-c_{it}\\
\log z_{it+1} & =\rho_{z}\log z_{it}+\psi_{it+1}\\
a_{it} & \geq0
\end{align*}
\end{itemize}
\end{frame}
%
\begin{frame}{Distributions and aggregates}
\begin{itemize}
\item <+->Household policy function $x^{*}$ where $x\in\{a,c,\ell\}$
function of:
\begin{itemize}
\item Individuals states $\left(z_{it},a_{it-1}\right)$
\item Current and future prices $\left\{ r_{\tau},w_{\tau}\right\} _{\tau\geq t}$
\end{itemize}
\item <+-> We can use those to compute the distribution of households $\boldsymbol{D}_{t}(z_{it},a_{it-1})$,
starting from $\boldsymbol{D}_{0}$ and applying the policies of periods
$0,\dots,t-1$
\item <+-> Aggregate policy are obtained combining policies and distributions:
\[
X^{hh}_{t}\left(\left\{ r_{\tau},w_{\tau}\right\} _{\tau\geq0}\right)=\int x^{\ast}_{t}(z_{it},a_{it-1},\left\{ r_{\tau},w_{\tau}\right\} _{\tau\geq t})d\boldsymbol{D}_{t}
\]
\item <+->When aggregating we \textbf{integrate} out individual states
\begin{itemize}
\item Aggregate $X^{hh}_{t}$ is only a function of the whole path $\left\{ r_{\tau},w_{\tau}\right\} _{\tau\geq0}$
(and of $\boldsymbol{D}_{0}$)
\end{itemize}
\item <+->$\Rightarrow$ If we know the path of prices $\left\{ r_{t},w_{t}\right\} $
we can calculate aggregate household behavior (consumption or savings)
\end{itemize}
\end{frame}
%
\begin{frame}{Equation system}

\vspace{-7mm} \small
\begin{align*}
\left[\begin{array}{c}
r^{K}_{t}-F_{K}(\Gamma_{t},K_{t-1},L_{t})\\
w_{t}-F_{L}(\Gamma_{t},K_{t-1},L_{t})\\
r_{t}-(r^{K}_{t}-\delta)\\
A_{t}-K_{t}\\
A_{t}-A^{hh}_{t}\\
L_{t}-L^{hh}_{t}\\
A^{hh}_{t}-\int a_{t}d\boldsymbol{D}_{t}\\
L^{hh}_{t}-\int\ell_{t}d\boldsymbol{D}_{t}\\
\underline{\boldsymbol{D}}_{t+1}-\Lambda^{\prime}_{t}\Pi^{\prime}_{z}\underline{\boldsymbol{D}}_{t}\\
a_{t}-a^{*}_{t}\\
\forall t\in\{0,1,\dots\},\text{given }\underline{\boldsymbol{D}}_{0}
\end{array}\right] & =\boldsymbol{0}
\end{align*}
\\
\vspace{-2mm} \normalsize
\begin{itemize}
\item <+->\textbf{Note}: Much larger system compared to Ramsey due to last
2 eqs.
\begin{itemize}
\item <+->$\boldsymbol{D}_{t},a^{*}_{t}$ define mass and optimal savings
policy at \textbf{the individual level}
\item <+->Standard Ramsey model: 8 eqs. per period
\item <+->HANC with $N_{z}=7,N_{a}=500:$ $8+2\times7\times500=7008$ per period
\end{itemize}
\end{itemize}
\end{frame}
%
\begin{frame}{Market clearing}
\begin{itemize}
\item \textbf{Capital market:} $K_{t}=A_{t}=\int a^{\ast}_{t}(z_{it},a_{it-1})d\boldsymbol{D}_{t}$
\item \textbf{Labor market:} $L_{t}=\int\ell^{\ast}_{t}(z_{it},a_{it-1})d\boldsymbol{D}_{t}=\int z_{it}d\boldsymbol{D}_{t}=1$
\item \textbf{Goods market:} $Y_{t}=\int c^{\ast}_{t}(z_{it},a_{it-1})d\boldsymbol{D}_{t}+I_{t}$\vspace{2mm}
\item \textbf{Walras:} Capital and labor market clears $\Rightarrow$ goods
market clears (using Euler's theorem)
\begin{align*}
C^{hh}_{t}+I_{t} & =\int c^{*}_{it}d\boldsymbol{D}_{t}+\left[K_{t}-(1-\delta)K_{t-1}\right]\\
 & =\int\left[(1+r_{t})a_{it-1}+w_{t}z_{it}-a_{it}\right]d\boldsymbol{D}_{t}\\
 & =\left[(1+r_{t})K_{t-1}+w_{t}L_{t}-K_{t}\right]+\left[K_{t}-(1-\delta)K_{t-1}\right]\\
 & =r^{K}_{t}K_{t-1}+w_{t}L_{t}\\
 & =Y_{t}
\end{align*}
\end{itemize}
\end{frame}
%

\section{Computing the Stationary Equilibrium}
\begin{frame}{Stationary equilibrium - equation system}

The \textbf{stationary equilibrium }satisfies
\begin{align*}
\left[\begin{array}{c}
r^{K}_{ss}-F_{K}(\Gamma_{ss},K_{ss},L_{ss})\\
w_{ss}-F_{L}(\Gamma_{ss},K_{ss},L_{ss})\\
r_{ss}-(r^{K}_{ss}-\delta)\\
A_{ss}-K_{ss}\\
A_{ss}-A^{hh}_{ss}\\
L_{ss}-L^{hh}_{ss}\\
A^{hh}_{ss}-\int a_{ss}d\boldsymbol{D}_{ss}\\
L^{hh}_{ss}-\int\ell_{ss}d\boldsymbol{D}_{ss}\\
\underline{\boldsymbol{D}}_{ss}-\Lambda^{\prime}_{ss}\Pi^{\prime}_{z}\underline{\boldsymbol{D}}_{ss}\\
a_{ss}-a^{*}_{ss}
\end{array}\right] & =\boldsymbol{0}
\end{align*}

\textbf{Note : }Households still move around >>inside<< the distribution
due to idiosyncratic shocks. Does not affect aggregates due to >>law
of large numbers<<
\end{frame}
%
\begin{frame}{Stationary equilibrium - more verbal definition}

{\small Given technology $\Gamma_{ss}$}{\small\par}
\begin{enumerate}
\item {\small Quantities $K_{ss}$ and $L_{ss},$ }{\small\par}
\item {\small prices $r_{ss}$ and $w_{ss}$ (always $\Pi_{ss}=0)$,}{\small\par}
\item {\small the distribution $\boldsymbol{D}_{ss}$ over $z_{it}$ and
$a_{it-1}$}{\small\par}
\item {\small and the policy functions $a^{\ast}_{ss}$, $\ell^{\ast}_{ss}$
and $c^{\ast}_{ss}$}{\small\par}
\end{enumerate}
{\small are such that}{\small\par}
\begin{enumerate}
\item {\small Households maximize expected utility (policy functions)}{\small\par}
\item {\small Firms maximize profits (prices)}{\small\par}
\item {\small$\boldsymbol{D}_{ss}$ is the invariant distribution implied
by the household problem}{\small\par}
\item {\small Mutual fund balance sheet is satisfied}{\small\par}
\item {\small The capital market clears}{\small\par}
\item {\small The labor market clears}{\small\par}
\item {\small The goods market clears}{\small\par}
\end{enumerate}
\end{frame}
%
\begin{frame}{How do we solve the household block in practice?}

The hard part is to solve the household block! How do we do it?
\begin{enumerate}
\item <+->Use EGM to obtain the policy functions $a_{ss},c_{ss}$ for a
given $r,w$
\begin{itemize}
\item The >>backward<< step
\end{itemize}
\item <+->Use the histogram method to obtain the stationary distribution
$\boldsymbol{D}_{ss}$
\begin{itemize}
\item The >>forward<< step
\end{itemize}
\item <+->Aggregate policy:
\[
A^{hh}_{ss}\left(\left\{ r_{ss},w_{ss}\right\} \right)=\int a^{\ast}_{ss}(z_{it},a_{it-1},\left\{ r_{ss},w_{ss}\right\} )d\boldsymbol{D}_{ss}
\]
\end{enumerate}

\end{frame}
%

\begin{frame}{Direct implementation (K guess)}

\textbf{Technology:} $F(K,L)=\Gamma K^{\alpha}L^{1-\alpha}$

\textbf{Root-finding problem} in $K_{ss}$ with the objective function:
\begin{enumerate}
\item Set $L_{ss}=1$ (and $\Pi_{ss}=0$)
\item Calculate $r_{ss}=\alpha\Gamma_{ss}(K_{ss})^{\alpha-1}-\delta$ and
$w_{ss}=(1-\alpha)\Gamma_{ss}(K_{ss})^{\alpha}$
\item Solve infinite horizon household problem \emph{backwards}, i.e. find
$\boldsymbol{a}^{\ast}_{ss}$
\item Simulate households \emph{forwards} until convergence, i.e. find $\boldsymbol{D}_{ss}$
\item Return $K_{ss}-\boldsymbol{a}^{\ast\prime}_{ss}\boldsymbol{D}_{ss}$\\
\end{enumerate}
\textbf{Note: $\boldsymbol{a}^{\ast\prime}_{ss}\boldsymbol{D}_{ss}=\sum_{i}a^{*}_{i,ss}D_{i}$}
\end{frame}
%
\begin{frame}{Direct implementation (r guess)}

\textbf{Technology:} $F(K,L)=\Gamma K^{\alpha}L^{1-\alpha}$

\textbf{Root-finding problem} in $r_{ss}$ with the objective function:
\begin{enumerate}
\item Set $L_{ss}=1$ (and $\Pi_{ss}=0$)
\item Calculate $K_{ss}=\left(\frac{r_{ss}+\delta}{\alpha\Gamma_{ss}}\right)^{\frac{1}{\alpha-1}}$
and $w_{ss}=(1-\alpha)\Gamma_{ss}(K_{ss})^{\alpha}$
\item Solve infinite horizon household problem \emph{backwards}, i.e. find
$\boldsymbol{a}^{\ast}_{ss}$
\item Simulate households \emph{forwards} until convergence, i.e. find $\boldsymbol{D}_{ss}$
\item Return $K_{ss}-\boldsymbol{a}^{\ast\prime}_{ss}\boldsymbol{D}_{ss}$
\end{enumerate}
\end{frame}

\begin{frame}{How to find a root numerically?}

  \begin{itemize}
    \item Finding a root of a function means finding $x^*$ such that $f(x^*)=0$
    \item How do we do it on a computer? Two main solutions (in the univariate case)
    \begin{itemize}
      \item[1.] Bisection
      \item[2.] Newton's method
    \end{itemize}
  \end{itemize}

\end{frame}


\begin{frame}{Root-finding: bisection}

\textbf{Equilibrium = root-finding problem: }find $x^{\ast}$ such that
$f(x^{\ast})=0$
\begin{itemize}
\item Here: $f(K)=A^{hh}_{ss}\left(r(K),w(K)\right)-K$
\item Each evaluation of $f$ is costly: a full solve of the household block
\end{itemize}
\vspace{2mm}

\onslide<2->\textbf{Bisection: }start from a bracket $[\underline{x},\overline{x}]$
with $f(\underline{x})f(\overline{x})<0$
\begin{enumerate}
\item <2->Compute the midpoint $m=(\underline{x}+\overline{x})/2$
\item <2->Stop if $|f(m)|<\text{tol}$
\item <2->Keep the half where the sign changes: if $f(\underline{x})f(m)<0$
set $\overline{x}=m$, else set $\underline{x}=m$
\end{enumerate}
\vspace{2mm}

\onslide<3->\textbf{Properties:}
\begin{itemize}
\item <3->Always converges if $f$ is continuous
\item <3->Slow: the bracket halves each iteration, error $\leq(\overline{x}-\underline{x})/2^{n}$
\item <3->Only works in one dimension
\end{itemize}
\end{frame}
%
\begin{frame}{Bisection: intuition}

\begin{center}
\begin{tikzpicture}[x=2.6cm,y=0.95cm,lbl/.style={font=\scriptsize},
bracket/.style={DarkRed,line width=1.6pt},
old/.style={black!35,line width=1.6pt}]

\path (0.3,-3.0) rectangle (3.9,3.3);
\draw[->] (0.4,0) -- (3.7,0) node[lbl,right]{$x$};
\draw[thick,domain=0.6:3.4,samples=50] plot (\x,{0.35*\x*\x-1});
\node[lbl,above left] at (2.8,1.744) {$f(x)$};

% brackets: current one on the axis, history stacked below
\only<1-2|handout:1>{\draw[bracket] (0.6,0) -- (3.4,0);
\node[lbl,below] at (0.6,0) {$\underline{x}$};
\node[lbl,below] at (3.4,0) {$\overline{x}$};
\fill (0.6,-0.874) circle (1.4pt) node[lbl,left]{$f(\underline{x})<0$};
\fill (3.4,3.046) circle (1.4pt) node[lbl,left]{$f(\overline{x})>0$};}
\only<3|handout:0>{\draw[bracket] (0.6,0) -- (2.0,0);
\node[lbl,below] at (0.6,0) {$\underline{x}$};
\node[lbl,below] at (2.0,0) {$\overline{x}$};}
\only<4|handout:0>{\draw[bracket] (1.3,0) -- (2.0,0);
\node[lbl,below] at (1.3,0) {$\underline{x}$};
\node[lbl,below] at (2.0,0) {$\overline{x}$};}
\only<5|handout:0>{\draw[bracket] (1.65,0) -- (2.0,0);
\node[lbl,below] at (1.65,0) {$\underline{x}$};
\node[lbl,below] at (2.0,0) {$\overline{x}$};}

% midpoints
\only<2|handout:1>{\draw[dashed] (2.0,0) -- (2.0,0.4);
\fill[DarkRed] (2.0,0.4) circle (1.6pt) node[lbl,above left,DarkRed]{$f(m)>0$};
\node[lbl,below,DarkRed] at (2.0,0) {$m$};}
\only<3|handout:0>{\draw[dashed] (1.3,0) -- (1.3,-0.4085);
\fill[DarkRed] (1.3,-0.4085) circle (1.6pt) node[lbl,below right,DarkRed]{$f(m)<0$};
\node[lbl,above,DarkRed] at (1.3,0) {$m$};}
\only<4|handout:0>{\draw[dashed] (1.65,0) -- (1.65,-0.047);
\fill[DarkRed] (1.65,-0.047) circle (1.6pt);
\node[lbl,above,DarkRed] at (1.65,0) {$m$};}

% history of brackets
\node[lbl,left] at (0.55,-1.5) {iteration 0};
\draw[alt=<1-2>{bracket}{old}] (0.6,-1.5) -- (3.4,-1.5);
\node[lbl,left,visible on=<3->] at (0.55,-1.9) {iteration 1};
\draw[alt=<3>{bracket}{old},visible on=<3->] (0.6,-1.9) -- (2.0,-1.9);
\node[lbl,left,visible on=<4->] at (0.55,-2.3) {iteration 2};
\draw[alt=<4>{bracket}{old},visible on=<4->] (1.3,-2.3) -- (2.0,-2.3);
\node[lbl,left,visible on=<5->] at (0.55,-2.7) {iteration 3};
\draw[alt=<5>{bracket}{old},visible on=<5->] (1.65,-2.7) -- (2.0,-2.7);

\end{tikzpicture}
\end{center}
\vspace{-1mm}

\begin{overlayarea}{\textwidth}{1.2cm}
\only<1|handout:0>{$f$ changes sign on $[\underline{x},\overline{x}]$, so
if $f$ is continuous there is a root inside.}
\only<2|handout:1>{$f(m)>0$, same sign as $f(\overline{x})$: the root
is to the left of $m$. Set $\overline{x}=m$.}
\only<3|handout:0>{$f(m)<0$, same sign as $f(\underline{x})$: the root
is to the right of $m$. Set $\underline{x}=m$.}
\only<4|handout:0>{Again $f(m)<0$: set $\underline{x}=m$.}
\only<5|handout:0>{Each evaluation of $f$ cuts the bracket in half, whatever
the shape of $f$.}
\end{overlayarea}
\end{frame}
%
\begin{frame}{Root-finding: Newton}

\begin{columns}[T]
\begin{column}{0.56\textwidth}
\textbf{Newton: }linearize $f$ around the guess $x_{n}$ and solve the
linear equation
\[
x_{n+1}=x_{n}-\frac{f(x_{n})}{f^{\prime}(x_{n})}
\]

\onslide<2->No analytical derivative? One extra evaluation of $f$:
$f^{\prime}(x)\approx\frac{f(x+h)-f(x)}{h}$
\end{column}
\begin{column}{0.44\textwidth}
\begin{center}
\begin{tikzpicture}[scale=0.95,lbl/.style={font=\scriptsize}]
\draw[->] (0.3,0) -- (3.5,0) node[lbl,right]{$x$};
\draw[thick,domain=0.6:3.2,samples=50] plot (\x,{0.35*\x*\x-1}) node[lbl,above]{$f(x)$};
\draw[DarkRed,thick] (3,2.15) -- (1.976,0);
\draw[dashed] (3,0) -- (3,2.15);
\draw[DarkRed,thick] (1.976,0.367) -- (1.711,0);
\draw[dashed] (1.976,0) -- (1.976,0.367);
\fill (3,2.15) circle (1.2pt);
\fill (1.976,0.367) circle (1.2pt);
\node[lbl,below] at (3,0) {$x_{0}$};
\node[lbl,below] at (2.05,0) {$x_{1}$};
\node[lbl,below] at (1.62,0) {$x_{2}$};
\end{tikzpicture}
\end{center}
\end{column}
\end{columns}
\vspace{-2mm}

\onslide<3->\textbf{Properties:}
\begin{itemize}
\item <3->Fast: very few iterations close to the root
\item <3->Not safe: can diverge or leave the domain from a bad guess
\item <3->Works in many dimensions (transition path, next lectures):\\
$\boldsymbol{x}_{n+1}=\boldsymbol{x}_{n}-\boldsymbol{J}(\boldsymbol{x}_{n})^{-1}\boldsymbol{f}(\boldsymbol{x}_{n})$
\end{itemize}

\end{frame}

\begin{frame}{}
\protect\phantomsection\label{section-1}
\begin{center}
\LARGE Let's code!
\end{center}
\end{frame}

\begin{frame}{Indirect implementation}

\textbf{Technology:} $F(K,L)=\Gamma K^{\alpha}L^{1-\alpha}$

\textbf{Consider $\Gamma_{ss}$ and $\delta$ as >>free<< parameters:}
\begin{enumerate}
\item Choose $r_{ss}$ and $w_{ss}$
\item Solve infinite horizon household problem \emph{backwards}, i.e. find
$\boldsymbol{a}^{\ast}_{ss}$
\item Simulate households \emph{forwards} until convergence, i.e. find $\boldsymbol{D}_{ss}$
\item Set $K_{ss}=\boldsymbol{a}^{\ast\prime}_{ss}\boldsymbol{D}_{ss}$
\item Set $L_{ss}=1$ (and $\Pi_{ss}=0$)
\item Set $\Gamma_{ss}=\frac{w_{ss}}{(1-\alpha)(K_{ss})^{\alpha}}$
\item Set $r^{K}_{ss}=\alpha\Gamma_{ss}(K_{ss})^{\alpha-1}$
\item Set $\delta=r^{k}_{ss}-r_{ss}$
\end{enumerate}
\end{frame}
%
\begin{frame}{Direct implementation ($\beta$)}

Set $r_{ss}=r^{target}$, $K_{ss}=K^{target}$, $Y_{ss}=Y^{target}$,
$w_{ss}L_{ss}/Y_{ss}=s^{target}_{L}$ and back out
\begin{enumerate}
\item $\alpha=1-s^{target}_{L}$
\item $\Gamma_{ss}=Y^{target}/(L^{1-\alpha}_{ss}K^{\alpha}_{ss})$
\item $\delta=\alpha Y^{target}/K^{target}-r^{target}$
\end{enumerate}
We know that $w_{ss}=(1-\alpha)Y^{target}.$ Then find the $\beta$
that clears the market

\textbf{Root-finding problem} in $\beta$ with the objective function:
\begin{enumerate}
\item Set $L_{ss}=1$ (and $\Pi_{ss}=0$),
\item Solve infinite horizon household problem \emph{backwards}, i.e. find
$\boldsymbol{a}^{\ast}_{ss}$for a given $\beta$
\item Simulate households \emph{forwards} until convergence, i.e. find $\boldsymbol{D}_{ss}$
\item Return $K_{ss}-\boldsymbol{a}^{\ast\prime}_{ss}\boldsymbol{D}_{ss}$
\end{enumerate}
\end{frame}

\begin{frame}{How to choose parameters?}
\begin{itemize}
\item <+->\textbf{External calibration: }Set subset of parameters to the
\emph{standard values in the literature} or \emph{directly from data
estimates} \\
(e.g. income process)
\item <+->\textbf{Internal calibration: }Set remaining parameters so the
model fit to a number of chosen \emph{macro-level and/or micro-level}
\emph{targets }based on empirical estimates
\begin{enumerate}
\item \textbf{Informal: }Roughly match targets by hand
\item \textbf{Formal: }

2a. Solve root-finding problem

2b. Minimize a squared loss function
\item \textbf{Estimation: }Formal with squared loss function (think GMM)
or likelihood function + standard errors
\end{enumerate}
\item <+->\textbf{Complication:} \emph{We must always solve for the steady
state for each guess of the parameters to be calibrated}
\end{itemize}
\end{frame}
%
\begin{frame}{Calibration at the quarterly level}
\begin{itemize}
\item $r_{ss}=0.05/4$ to match 5\% annual interest rate
\item $Y_{ss}=1$ as a normalization
\item $K_{ss}=16$ to match annual wealth-to-output ratio of 4
\item $\alpha=1/3$ to match labor share of roughly 2/3
\item $\sigma_{\psi}=0.5$, $\rho=0.9$: data on income inequality and risk
\end{itemize}
\end{frame}
%



\section{Some Properties of the HANC steady-state}

\begin{frame}{Consumption function}
\begin{itemize}
\item \textbf{Euler-equation still necessary for $a_{it}>0$:
\[
c^{-\sigma}_{it}=\beta(1+r_{t+1})\mathbb{E}_{t}\left[c^{-\sigma}_{it+1}\right]
\]
}
\item \textbf{Precautionary saving:}
\begin{enumerate}
\item Low consumption for low cash-on-hand $\rightarrow$ \emph{buffer-stock
target}
\item Steep slope for low cash-on-hand $\rightarrow$ \emph{high MPC}
\end{enumerate}
\end{itemize}
\includegraphics[width=1\textwidth]{figs/c_func}
\end{frame}
%
\begin{frame}{Some amount of inequality}
\begin{itemize}
\item \textbf{Productivity: }Marginal distribution over only $z_{it}$
\item \textbf{Savings: }Marginal distribution over only $a_{it}$
\end{itemize}
\includegraphics[width=1\textwidth]{figs/distribution}
\begin{itemize}
\item \textbf{Drivers of wealth inequality here: income shocks}
\end{itemize}
\end{frame}
%
\begin{frame}{Steady state interest rate}
\begin{itemize}
\item <+->\textbf{Representative agent / complete markets: }

Derived from aggregate Euler-equation
\[
C^{-\sigma}_{t}=\beta(1+r_{t+1})C^{-\sigma}_{t+1}\Rightarrow C^{-\sigma}_{ss}=\beta(1+r_{ss})C^{-\sigma}_{ss}\Leftrightarrow\beta=\frac{1}{1+r_{ss}}
\]

\item <+->\textbf{Heterogeneous agents: }\emph{No such equation exists}
\end{itemize}
\begin{figure}
\centering\includegraphics[width=0.6\linewidth]{figs/market_clearing}
\end{figure}
\end{frame}
%
\begin{frame}{Risk drives wealth accumulation}

\includegraphics[width=1\textwidth]{figs/sigma_psi_vs_A}

\end{frame}
%
\begin{frame}{Marginal Propensity to Consume}

\includegraphics[width=1\textwidth]{figs/MPC_policies}
\end{frame}
%
\begin{frame}{Tradeoff between matching aggregate wealth and MPCs}

\includegraphics[width=1\textwidth]{figs/MPC_vs_K}

How to solve this? One solution: add ex-ante heterogeneity in discount
factors.
\end{frame}
%

\section{Exercises}

\begin{frame}{Exercise 1: HANC with ex-ante heterogeneity}

Add permanent $\beta$ heterogeneity to the HANC model:

\begin{align*}
v_{t}(\beta_{i},z_{it},a_{it-1}) & =\max_{c_{it}}\frac{c^{1-\sigma}_{it}}{1-\sigma}+\beta_{i}\mathbb{E}_{t}\left[v_{it+1}(z_{it+1},a_{it})\right]\\
\text{s.t. }a_{it}+c_{it} & =(1+r_{t})a_{it-1}+w_{t}z_{it},\,\,\,a_{it}\geq0\\
\log z_{it+1} & =\rho_{z}\log z_{it}+\psi_{it+1}\,\,\,,\psi_{it}\sim\mathcal{N}(\mu_{\psi},\sigma_{\psi}),\,\mathbb{E}[z_{it}]=1\,
\end{align*}

Assume that we have three types of households: $\beta_{i}\in\{\beta-\delta^{\beta},\beta,\beta+\delta^{\beta}\}.$
Find $\delta^{\beta}$ such that $MPC=0.27$ and $K/Y=16$.
\end{frame}
%
\begin{frame}{Exercise 2: HANCGovModel}
\begin{itemize}
\item \textbf{No production. }No physical savings instrument
\item \textbf{Households: }Get stochastic endowment\textbf{ }$z_{it}$ of
consumption good
\item \textbf{Government:}
\begin{enumerate}
\item Choose government spending
\item Collect taxes, $\tau_{t}$, proportional to endowment
\item Bonds: Pays $1$ unit of the consumption good next period. Price is
$p^{B}_{t}<1$
\begin{align*}
p^{B}_{t}B_{t} & +\int\tau_{t}z_{it}d\boldsymbol{D}_{t}=B_{t-1}+G_{t}\\
\tau_{t} & =\tau_{ss}+\varphi\left(B_{t-1}-B_{ss}\right)
\end{align*}
\end{enumerate}
\item \textbf{Market clearing:
\begin{align*}
B_{t} & =A^{hh}_{t}\\
C^{hh}_{t}+G_{t} & =\int z_{it}d\boldsymbol{D}_{t}=1
\end{align*}
}
\end{itemize}
\end{frame}
%
\begin{frame}{Exercise 2: Households}

\textbf{Households:}
\begin{align*}
v_{t}(z_{it},a_{it-1}) & =\max_{c_{it}}\frac{c^{1-\sigma}_{it}}{1-\sigma}+\beta\mathbb{E}_{t}\left[v_{it+1}(z_{it+1},a_{it})\right]\\
\text{s.t. }p^{B}_{t}a_{it}+c_{it} & =a_{it-1}+(1-\tau_{t})z_{it},\,\,\,a_{it}\geq0\\
\log z_{it+1} & =\rho_{z}\log z_{it}+\psi_{it+1}\,\,\,,\psi_{it}\sim\mathcal{N}(\mu_{\psi},\sigma_{\psi}),\,\mathbb{E}[z_{it}]=1\,
\end{align*}

\textbf{Euler-equation}:
\begin{align*}
c^{-\sigma}_{it} & =\beta\frac{\underline{v}_{a,t+1}(z_{it},a_{it})}{p^{B}_{t}}
\end{align*}

\textbf{Envelope condition:
\[
\underline{v}_{a,t}(z_{it-1},a_{it-1})=c^{-\sigma}_{it}
\]
}
\end{frame}
%
\begin{frame}{Exercise 2: Questions}
\begin{enumerate}
\item \textbf{Define the stationary equilibrium}
\item \textbf{Solve and simulate the household problem }\\
with $p^{B}_{ss}=0.975$ and $\tau_{ss}=0.12$.
\item \textbf{Find the stationary equilibrium}\\
with $G_{ss}=0.10$ and $\tau_{ss}=0.12$.
\item \textbf{What happens for $\tau_{ss}\in(0.11,0.15)?$ }
\item \textbf{When is average household utility maximized? }
\end{enumerate}
\vspace{5mm}\textbf{Note:} Full solution available at this \href{https://github.com/NumEconCopenhagen/GEModelToolsNotebooks/tree/master/HANC-gov}{\underline{link}}
\end{frame}
%

\end{document}
